% Options for packages loaded elsewhere
\PassOptionsToPackage{unicode}{hyperref}
\PassOptionsToPackage{hyphens}{url}
\documentclass[
  11pt,
]{article}
\usepackage{xcolor}
\usepackage[left=3cm,right=3cm,top=2.5cm,bottom=2.5cm]{geometry}
\usepackage{amsmath,amssymb}
\setcounter{secnumdepth}{5}
\usepackage{iftex}
\ifPDFTeX
  \usepackage[T1]{fontenc}
  \usepackage[utf8]{inputenc}
  \usepackage{textcomp} % provide euro and other symbols
\else % if luatex or xetex
  \usepackage{unicode-math} % this also loads fontspec
  \defaultfontfeatures{Scale=MatchLowercase}
  \defaultfontfeatures[\rmfamily]{Ligatures=TeX,Scale=1}
\fi
\usepackage{lmodern}
\ifPDFTeX\else
  % xetex/luatex font selection
\fi
% Use upquote if available, for straight quotes in verbatim environments
\IfFileExists{upquote.sty}{\usepackage{upquote}}{}
\IfFileExists{microtype.sty}{% use microtype if available
  \usepackage[]{microtype}
  \UseMicrotypeSet[protrusion]{basicmath} % disable protrusion for tt fonts
}{}
\usepackage{setspace}
\makeatletter
\@ifundefined{KOMAClassName}{% if non-KOMA class
  \IfFileExists{parskip.sty}{%
    \usepackage{parskip}
  }{% else
    \setlength{\parindent}{0pt}
    \setlength{\parskip}{6pt plus 2pt minus 1pt}}
}{% if KOMA class
  \KOMAoptions{parskip=half}}
\makeatother
\usepackage{longtable,booktabs,array}
\usepackage{caption}
\captionsetup[table]{skip=6pt}
\newcounter{none} % for unnumbered tables
\usepackage{calc} % for calculating minipage widths
% Correct order of tables after \paragraph or \subparagraph
\usepackage{etoolbox}
\makeatletter
\patchcmd\longtable{\par}{\if@noskipsec\mbox{}\fi\par}{}{}
\makeatother
% Allow footnotes in longtable head/foot
\IfFileExists{footnotehyper.sty}{\usepackage{footnotehyper}}{\usepackage{footnote}}
\makesavenoteenv{longtable}
\usepackage{graphicx}
\makeatletter
\newsavebox\pandoc@box
\newcommand*\pandocbounded[1]{% scales image to fit in text height/width
  \sbox\pandoc@box{#1}%
  \Gscale@div\@tempa{\textheight}{\dimexpr\ht\pandoc@box+\dp\pandoc@box\relax}%
  \Gscale@div\@tempb{\linewidth}{\wd\pandoc@box}%
  \ifdim\@tempb\p@<\@tempa\p@\let\@tempa\@tempb\fi% select the smaller of both
  \ifdim\@tempa\p@<\p@\scalebox{\@tempa}{\usebox\pandoc@box}%
  \else\usebox{\pandoc@box}%
  \fi%
}
% Set default figure placement to htbp
\def\fps@figure{htbp}
\makeatother
% definitions for citeproc citations
\NewDocumentCommand\citeproctext{}{}
\NewDocumentCommand\citeproc{mm}{%
  \begingroup\def\citeproctext{#2}\cite{#1}\endgroup}
\makeatletter
 % allow citations to break across lines
 \let\@cite@ofmt\@firstofone
 % avoid brackets around text for \cite:
 \def\@biblabel#1{}
 \def\@cite#1#2{{#1\if@tempswa , #2\fi}}
\makeatother
\newlength{\cslhangindent}
\setlength{\cslhangindent}{1.5em}
\newlength{\csllabelwidth}
\setlength{\csllabelwidth}{3em}
\newenvironment{CSLReferences}[2] % #1 hanging-indent, #2 entry-spacing
 {\begin{list}{}{%
  \setlength{\itemindent}{0pt}
  \setlength{\leftmargin}{0pt}
  \setlength{\parsep}{0pt}
  % turn on hanging indent if param 1 is 1
  \ifodd #1
   \setlength{\leftmargin}{\cslhangindent}
   \setlength{\itemindent}{-1\cslhangindent}
  \fi
  % set entry spacing
  \setlength{\itemsep}{#2\baselineskip}}}
 {\end{list}}
\usepackage{calc}
\newcommand{\CSLBlock}[1]{\hfill\break\parbox[t]{\linewidth}{\strut\ignorespaces#1\strut}}
\newcommand{\CSLLeftMargin}[1]{\parbox[t]{\csllabelwidth}{\strut#1\strut}}
\newcommand{\CSLRightInline}[1]{\parbox[t]{\linewidth - \csllabelwidth}{\strut#1\strut}}
\newcommand{\CSLIndent}[1]{\hspace{\cslhangindent}#1}
\setlength{\emergencystretch}{3em} % prevent overfull lines
\providecommand{\tightlist}{%
  \setlength{\itemsep}{0pt}\setlength{\parskip}{0pt}}
% Paragraph scaffolding for manuscript revision.
%
% Two environments, presented in this order per paragraph:
%
%   rgt   -- author's rewrite (initially a placeholder)
%            small, italic, dark blue
%   orig  -- original Claude-authored draft, and until the rewrite
%            replaces it, the body text of the manuscript
%            normal size, gray
%
% The `bullets' environment (a boiled-down topic sentence plus bullet
% points, one per paragraph) is no longer used inside manuscripts:
% those blocks now live in a companion bullets.Rmd supplement in the
% same report directory, which typesets them as plain \textbf{} +
% itemize and needs no part of this file. The definition is retained
% only so that a stray block left in a document still compiles.
%
% This file is identical across every manuscript that includes it.
% It previously existed in two variants that swapped the sizes of
% rgt and orig, so the same markup rendered differently depending on
% which repository the manuscript lived in. The variant retained here
% keeps orig at full size because orig, not rgt, currently carries
% the prose in every manuscript in the pool.

\usepackage{xcolor}

\newcounter{pgraph}

\newenvironment{bullets}{%
  \stepcounter{pgraph}%
  \begin{quote}\small\sffamily\color[gray]{0.30}%
  {\normalfont\scriptsize\color[gray]{0.58}[\thepgraph]\enspace}}{%
  \end{quote}}

\newenvironment{rgt}{%
  \begin{quote}\small\itshape\color[RGB]{20,60,160}}{%
  \end{quote}}

\newenvironment{orig}{%
  \begin{quote}\normalsize\color[gray]{0.55}}{%
  \end{quote}}

% backward compat alias
\newenvironment{claudecode}{\begin{orig}}{\end{orig}}
%% tools/stamp.tex
%% Provenance footer for PDFs rendered from this project.
%% Vendored by zzcollab; this file is not project-specific. The
%% footer prints the source document and the time of rendering.
%% The git version is defined here too, and so is available to a
%% project that wants it on the page, but it is left off the
%% default footer: it already names the staged copy in share/ and
%% appears in its manifest row. All values are supplied at render
%% time by tools/stamp-render.R, which writes a small generated
%% file that redefines the macros below. The \providecommand
%% defaults keep a bare LaTeX compile working when the document is
%% built without the wrapper.
\usepackage{fancyhdr}
\usepackage{xcolor}
\providecommand{\stampsource}{\detokenize{(source not recorded)}}
\providecommand{\stamptime}{(time not recorded)}
\providecommand{\stampversion}{\detokenize{(version not recorded)}}
\setlength{\footskip}{40pt}
\newcommand{\stampfootcontent}{%
  \footnotesize\color{gray}%
  \parbox[b]{\textwidth}{%
    \stamptime\hfill\thepage\\[2pt]%
    \ttfamily\raggedright\stampsource}}
\pagestyle{fancy}
\fancyhf{}
\fancyfoot[C]{\stampfootcontent}
\renewcommand{\headrulewidth}{0pt}
\renewcommand{\footrulewidth}{0.4pt}
%% The title page uses the `plain` page style; redefine it so the
%% stamp appears there too.
\fancypagestyle{plain}{%
  \fancyhf{}%
  \fancyfoot[C]{\stampfootcontent}%
  \renewcommand{\headrulewidth}{0pt}%
  \renewcommand{\footrulewidth}{0.4pt}}
\renewcommand{\stampsource}{\textasciitilde{}/\allowbreak{}prj/\allowbreak{}res/\allowbreak{}36-pmsimstats-ng/\allowbreak{}pmsimstats-ng/\allowbreak{}analysis/\allowbreak{}report/\allowbreak{}11-combined-dgp-architecture/\allowbreak{}report.Rmd}
\renewcommand{\stamptime}{2026-08-15 16:07 PDT}
\renewcommand{\stampversion}{\detokenize{bc0b671-wip}}
\usepackage{amsmath}
\usepackage{amssymb}
\usepackage{booktabs}
\usepackage{longtable}
\usepackage{float}
\usepackage[mathlines]{lineno}
\linenumbers
\usepackage{bookmark}
\IfFileExists{xurl.sty}{\usepackage{xurl}}{} % add URL line breaks if available
\urlstyle{same}
% fallback for those not using the hyperref driver hyperxmp:
\makeatletter
\@ifundefined{xmpquote}{\newcommand{\xmpquote}[1]{#1}}{}
\makeatother
\hypersetup{
  pdftitle={A combined data-generating architecture for biomarker-treatment interactions: channel super-additivity and mixing-weight sensitivity in aggregated N-of-1 trials},
  pdfauthor={pmsimstats team},
  hidelinks,
  pdfcreator={LaTeX via pandoc}}

\title{A combined data-generating architecture for biomarker-treatment interactions: channel super-additivity and mixing-weight sensitivity in aggregated N-of-1 trials}
\author{pmsimstats team}
\date{August 15, 2026}

\begin{document}
\maketitle

{
\setcounter{tocdepth}{2}
\tableofcontents
}
\setstretch{1.1}
\section{Abstract}

\textbf{Background.} A companion paper\textsuperscript{\citeproc{ref-pmsimstats-paper01}{1}} distinguishes
two data-generating architectures for biomarker-treatment
interactions in aggregated N-of-1 trials: direct mean moderation
(the mean-moderation architecture), in which the biomarker scales the treatment
effect in the conditional mean, and differential correlation
(the covariance-moderation architecture, the multivariate-normal approach), in which the
interaction is encoded as treatment-state-dependent correlation in
the joint distribution. That paper treats A and B as alternatives.
Here we consider the case in which both channels are active at once.

\textbf{Methods.} We define the dual-channel architecture, which
carries independent parameters \(c_{bm,A}\) (mean channel) and
\(c_{bm,B}\) (covariance channel), and which reduces algebraically to
the pure mean-moderation architecture when \(c_{bm,B} = 0\) and to pure
the covariance-moderation architecture when \(c_{bm,A} = 0\). We show analytically that the
biomarker-treatment signal delivered to the fixed analysis model
is additive across the two channels, and that, because the
Monte-Carlo power function is convex in the noncentrality over the
sub-ceiling regime, the two channels combine super-additively in
power. We then quantify the combined architecture by Monte Carlo
simulation on a \(3 \times 3\) grid of channel weights crossed with
three within-subject designs and three carryover half-lives.

\textbf{Results.} {[}To be finalized after the common-driver boundary
re-run; see Section 4.2.{]} Empirical power increases monotonically in
each channel weight, and interior cells exceed the probability-scale
additive benchmark, confirming the analytical prediction of
super-additivity. The covariance and mean channels contribute
comparably to carryover sensitivity, and the design ranking
established for the pure architectures is preserved throughout the
combined grid.

\textbf{Conclusions.} A combined architecture is the appropriate
data-generating model when a biomarker is hypothesized to act
through both a pharmacodynamic-scaling pathway and a co-regulatory
coupling pathway. It should not be a default: it carries a second
free parameter and requires independent evidence for each channel.
Its principal use is as a sensitivity framework spanning the two
pure-architecture extremes.

\bigskip\noindent

\rule{\textwidth}{0.4pt}\bigskip

\section{1. Introduction}

\begin{rgt}
To be completed by rgt.

\end{rgt}

\begin{orig}
A companion paper\textsuperscript{\citeproc{ref-pmsimstats-paper01}{1}} formalizes two data-generating
process (DGP) architectures for simulation studies of predictive
biomarker-treatment interactions in aggregated N-of-1 trials, and
shows that the choice between them materially changes the estimated
power loss under carryover. Under the mean-moderation architecture the biomarker
scales the treatment effect in the conditional mean; under
the covariance-moderation architecture the interaction is carried by
treatment-state-dependent correlation between biomarker and response
in the joint distribution. That paper treats the two as competing
choices. The present paper takes up the case, deferred there, in
which both mechanisms operate at once. We call the resulting
data-generating model the dual-channel architecture. It is not a third
alternative on the same footing as A and B; it is their common
generalization, with A and B recovered exactly as its
single-channel limits.

\end{orig}

\begin{rgt}
To be completed by rgt.

\end{rgt}

\begin{orig}
The paper makes two contributions. First, we give an analytical
account (Section\textasciitilde3) of why the two channels, when active together,
combine super-additively in power rather than merely adding. The
argument identifies the additivity of the biomarker-treatment
signal delivered to the analysis model as the mechanism, and the
convexity of the Monte-Carlo power function over the sub-ceiling
regime as the reason additivity of signal becomes super-additivity
of power. This addresses a gap noted for the descriptive combined
results reported in the companion paper, where the super-additivity
was observed but not explained. Second, we characterize the combined
architecture empirically over a \(3 \times 3\) grid of channel weights
crossed with three within-subject designs and three carryover
half-lives, with the boundary cells validated numerically against
the pure-architecture baselines of the companion paper.

\end{orig}

\bigskip\noindent

\rule{\textwidth}{0.4pt}\bigskip

\section{2. The Combined Architecture}

\subsection{2.1 Definition}

\begin{rgt}
To be completed by rgt.

\end{rgt}

\begin{orig}
the dual-channel architecture applies both interaction mechanisms in
sequence. First, the multivariate-normal covariance matrix is
constructed with biomarker-to-biological-response (BM-BR)
differential correlation weighted by \(c_{bm,B}\), exactly as under
the covariance-moderation architecture: the on-drug correlation is \(c_{bm,B}\), decaying to
\(c_{bm,B} \cdot e^{-\lambda t_{sd}}\) during off-drug periods, where
\(t_{sd}\) is time since discontinuation. Second, after the MVN draw,
the BR component is additively shifted at on-drug timepoints by
\[
\Delta BR_{it} = c_{bm,A} \cdot z_{B,i} \cdot \sigma_{BR}
  \cdot \mathbf{1}[\text{on-drug}],
\]
where \(z_{B,i} = (B_i - \mu_B)/\sigma_B\) is the standardized
biomarker for participant \(i\), exactly as under the mean-moderation architecture.

\end{orig}

\subsection{2.2 Boundary reductions and identifiability}

\begin{rgt}
To be completed by rgt.

\end{rgt}

\begin{orig}
The two parameters \(c_{bm,A}\) and \(c_{bm,B}\) are independent.
Setting \(c_{bm,B} = 0\) removes the covariance channel and recovers
the pure mean-moderation architecture; setting \(c_{bm,A} = 0\) removes the mean-shift
channel and recovers the pure covariance-moderation architecture. This combined form is the
mean-and-covariance-structure model of Arminger et al.\textsuperscript{\citeproc{ref-arminger1999}{2}}, the general psychometric parameterization within
which the mean-moderation and covariance-moderation architectures are constrained submodels; a companion
paper\textsuperscript{\citeproc{ref-pmsimstats-paper03}{3}} situates it as one node of a
biomarker-moderation spectrum bridging continuous-heterogeneity and
latent-class encodings. The boundary cases are
enforced algebraically rather than by convention, so they serve as
internal validation gates for any simulation study using
the dual-channel architecture: a correctly implemented combined driver must
reproduce the pure-architecture power values at the boundaries when
run under matched design, sample-size, and calibration settings.
Section\textasciitilde4.2 uses exactly this property to validate the simulation
against the companion paper's baselines.

We note that when both parameters are set to their full values
(\(c_{bm,A} = c_{bm,B} = 0.45\)), the total interaction effect on the
BR mean is approximately double that of either pure architecture at
\(c_{bm} = 0.45\). A researcher wishing to hold total effect size
constant across architectures may instead parameterize as
\(c_{bm,A} = (1-\alpha) \cdot c\) and \(c_{bm,B} = \alpha \cdot c\)
for a mixing weight \(\alpha \in [0,1]\) and total weight \(c\). We do
not adopt this constraint in the simulations below, preferring the
factorial parameterization that spans the full
\((c_{bm,A}, c_{bm,B})\) space.

\end{orig}

\subsection{2.3 Biological rationale}

\begin{rgt}
To be completed by rgt.

\end{rgt}

\begin{orig}
The two channels are mechanistically distinct and need not scale in
proportion. The mean-shift channel (the mean-moderation architecture) reflects a
pharmacodynamic relationship: the drug's biological effect on BR is
amplified in participants whose biomarker is above the population
mean, because the biomarker governs the effective dose or receptor
occupancy. This mechanism operates at the level of the mean
trajectory for each participant. The covariance channel
(the covariance-moderation architecture) reflects a different process: on-drug periods
impose a co-regulatory coupling between biomarker expression and
treatment response, perhaps because the drug activates a pathway
whose output is gated by the biomarker signal. This coupling is
absent during off-drug periods, producing the differential
correlation structure that the covariance-moderation architecture encodes in the covariance
matrix. A given drug could activate one, both, or neither channel.
Cast in latent-class terms, both channels carry information about an
unobserved responder class: the mean channel through the
class-conditional expectation \(E[BR \mid B]\) and the covariance
channel through the class-conditional dispersion
\(\mathrm{Var}[BR \mid B]\)\textsuperscript{\citeproc{ref-pmsimstats-paper03}{3}}. The linear-mixed
analogue, in which the heterogeneity resides in the random-effects
population, is developed by Verbeke and Lesaffre\textsuperscript{\citeproc{ref-verbekelesaffre1996}{4}}.

\end{orig}

\bigskip\noindent

\rule{\textwidth}{0.4pt}\bigskip

\section{3. Channel Super-additivity: an Analytical Account}

\begin{rgt}
To be completed by rgt.

\end{rgt}

\begin{orig}
The fixed analysis model detects the interaction through the
coefficient on the biomarker-by-exposure term, whose estimate is
driven by the on-drug covariance between the biomarker \(B\) and the
biological response \(BR\). Consider that covariance under each
channel. Under the covariance channel alone (the covariance-moderation architecture), the
on-drug joint distribution sets
\[
\mathrm{Cov}(B,\, BR \mid \text{on-drug})
  = c_{bm,B}\, \sigma_B \sigma_{BR}.
\]
Under the mean channel alone (the mean-moderation architecture), the on-drug shift
\(\Delta BR = c_{bm,A}\, z_B\, \sigma_{BR}
 = c_{bm,A}\, (B - \mu_B)\, \sigma_{BR} / \sigma_B\) contributes
\[
\mathrm{Cov}(B,\, \Delta BR)
  = \frac{c_{bm,A}\, \sigma_{BR}}{\sigma_B}\, \mathrm{Var}(B)
  = c_{bm,A}\, \sigma_B \sigma_{BR}.
\]
Under the combined architecture the on-drug response is the sum of
the correlated draw and the mean shift, so the two covariance
contributions add:
\[
\mathrm{Cov}(B,\, BR_C \mid \text{on-drug})
  = (c_{bm,A} + c_{bm,B})\, \sigma_B \sigma_{BR}.
\]
The biomarker-treatment signal delivered to the analysis model
is therefore exactly additive in the two channel weights. To
leading order the residual variance is unchanged by the channels
(the covariance channel preserves the marginal BR variance, and the
mean channel inflates it only by the second-order term
\(c_{bm,A}^2 \sigma_{BR}^2\)), so the noncentrality parameter of the
interaction test is, to first order, additive across channels:
\(\eta_C \approx \eta_A + \eta_B\).

\end{orig}

\begin{rgt}
To be completed by rgt.

\end{rgt}

\begin{orig}
Additivity of the noncentrality does not imply additivity of power,
because power is a nonlinear function of the noncentrality. Write
\(\pi(\eta)\) for the rejection probability of the interaction test as
a function of its noncentrality \(\eta \ge 0\). On the range over
which \(\pi\) is below approximately one half, \(\pi\) is convex in
\(\eta\): each additional unit of noncentrality buys more power than
the last, until the function inflects and saturates toward one. For
\(\eta_A, \eta_B\) in this convex regime, convexity gives
\[
\pi(\eta_A + \eta_B) - \pi(0) >
 \bigl[\pi(\eta_A) - \pi(0)\bigr] +
 \bigl[\pi(\eta_B) - \pi(0)\bigr],
\]
which is precisely super-additivity relative to the
probability-scale additive benchmark
\(p(c_{bm,A}, 0) + p(0, c_{bm,B}) - p(0,0)\). The combined power
exceeds the benchmark exactly when the two single-channel cells sit
in the convex, sub-ceiling part of the power curve, and the excess
shrinks, and can reverse to sub-additivity, once cells approach the
ceiling and enter the concave regime. The empirical grid of
Section\textasciitilde4 lies in the sub-ceiling regime (single-channel powers
between roughly 0.03 and 0.30), so super-additivity is the predicted
and observed pattern there.

\end{orig}

\bigskip\noindent

\rule{\textwidth}{0.4pt}\bigskip

\section{4. Simulation Study}

\subsection{4.1 Design}

\begin{rgt}
To be completed by rgt.

\end{rgt}

\begin{orig}
We ran a factorial simulation on the dual-channel architecture parameter grid.
The design factors are the crossover (CO), Hybrid, and
open-label-with-blinded-discontinuation (OL+BDC) designs;
\(t_{1/2} \in \{0, 0.5, 1.0\}\) weeks; 1000 replicates per cell; and
the E2 analysis specification. The channel-weight panel is
\(c_{bm,A} \times c_{bm,B} \in \{0, 0.22, 0.45\}^2\), nine cells per
design-carryover combination. The analysis model is held fixed
throughout and is identical to that of the companion paper: an
\texttt{nlme::lme} fit with \texttt{corCAR1} continuous-time
residual correlation and a single biomarker-by-exposure interaction
term.

\end{orig}

\subsection{4.2 Boundary validation against the companion baselines}

\begin{rgt}
To be completed by rgt.

\end{rgt}

\begin{orig}
\textbf{Status: pending re-run.} The combined-architecture results
must be produced by a driver whose sample-size convention and
Gompertz and biomarker calibration match the companion paper's
production driver exactly, so that the boundary cells reproduce the
pure-architecture power values numerically and not merely
algebraically. An earlier combined driver used a different
sample-size convention (\(N = 35\) as a total rather than per
randomization path), which made its boundary cells uniformly lower
than, and not directly comparable to, the companion baselines. This
section will report a boundary-validation table placing the combined
driver's \(c_{bm,B} = 0\) column and \(c_{bm,A} = 0\) row alongside the
companion paper's the mean-moderation architecture and the covariance-moderation architecture baselines, and
will confirm agreement to within Monte Carlo standard error
(approximately 0.013 to 0.016 at these cell sizes). The results in
Section\textasciitilde4.3 are reported at the total-sample convention of the
earlier driver and will be re-stamped once the matched re-run is
complete.

\end{orig}

\subsection{4.3 Empirical power over the channel-weight grid}

\begin{table}[h]
\small\centering
\caption{the dual-channel architecture empirical power: $3 \times 3$ channel-weight
panel at carryover $t_{1/2} = 1.0$ week, E2 analysis specification,
1000 replicates per cell. Rows index $c_{bm,A}$ (mean channel);
columns index $c_{bm,B}$ (covariance channel). Boundary column
$c_{bm,B} = 0$ corresponds to the pure mean-moderation architecture; boundary row
$c_{bm,A} = 0$ corresponds to the pure covariance-moderation architecture. Values are at the
earlier driver's total-sample convention and are pending the matched
re-run of Section 4.2; Monte Carlo standard error columns will be
added at that time.}
\label{tab:arch-c-power}
\begin{tabular}{lrrr}
\toprule
& \multicolumn{3}{c}{$c_{bm,B}$ (covariance channel)} \\
\cmidrule(l){2-4}
$c_{bm,A}$ (mean channel) & 0.00 & 0.22 & 0.45 \\
\midrule
\multicolumn{4}{l}{\textit{Crossover (CO)}} \\
0.00 & 0.046 & 0.064 & 0.206 \\
0.22 & 0.070 & 0.179 & 0.361 \\
0.45 & 0.172 & 0.354 & 0.528 \\
\addlinespace
\multicolumn{4}{l}{\textit{Hybrid}} \\
0.00 & 0.034 & 0.083 & 0.273 \\
0.22 & 0.108 & 0.299 & 0.621 \\
0.45 & 0.389 & 0.642 & 0.870 \\
\addlinespace
\multicolumn{4}{l}{\textit{OL+BDC}} \\
0.00 & 0.017 & 0.049 & 0.148 \\
0.22 & 0.086 & 0.164 & 0.362 \\
0.45 & 0.276 & 0.473 & 0.688 \\
\bottomrule
\end{tabular}
\end{table}

\begin{rgt}
To be completed by rgt.

\end{rgt}

\begin{orig}
Table\textasciitilde{}\ref{tab:arch-c-power} reports empirical power across the full
\(3 \times 3\) channel-weight panel at the most demanding carryover
condition (\(t_{1/2} = 1.0\) week). Three features merit comment.
First, the monotone structure is exact: power increases strictly
with each channel weight, holding the other fixed, across all three
designs, confirming that the channels contribute independently to
detectability. Second, interior cells exceed the probability-scale
additive benchmark. At (\(c_{bm,A} = 0.22\), \(c_{bm,B} = 0.22\)) in the
Hybrid design, power is 0.299; the additive baseline
\(p(0.22,0) + p(0,0.22) - p(0,0) = 0.108 + 0.083 - 0.034 = 0.157\) is
exceeded by 0.142, the super-additivity predicted in Section\textasciitilde3.
Third, the covariance and mean channels contribute comparably to
carryover sensitivity, and design choice remains consequential: at
the full-signal, one-week-carryover cell, Hybrid (0.870)
substantially outperforms OL+BDC (0.688) and CO (0.528).

\end{orig}

\bigskip\noindent

\rule{\textwidth}{0.4pt}\bigskip

\section{5. When a Combined Architecture is Appropriate}

\begin{rgt}
To be completed by rgt.

\end{rgt}

\begin{orig}
the dual-channel architecture is appropriate when there is independent evidence
that both the pharmacodynamic-scaling pathway and the co-regulatory
coupling pathway are biologically active for the same
biomarker-drug pair. Such evidence might include mechanistic assay
data (for example, the drug enhances biomarker-outcome correlation
in vitro, consistent with the B channel, while receptor occupancy
scales with biomarker level, consistent with the A channel) or
replication across independent trials individually underpowered to
distinguish the channels. The dual-channel architecture should not be the default
choice for simulation studies despite its generality. It carries two
free parameters where either pure architecture requires one, and
calibrating both independently demands more prior information than
is typically available at the trial-design stage. This estimation
caution is not a claim that the combined mechanism is biologically
unusual: a companion paper\textsuperscript{\citeproc{ref-pmsimstats-paper03}{3}} argues that the combined
mean-and-covariance form is plausibly the most biologically faithful
regime for biomarker-guided psychiatric trials, because a biomarker
that indexes a responder subtype will in general shift both the
class-conditional mean and the class-conditional dispersion of
response. The recommendation against the dual-channel architecture as a default is
therefore about identifiability and prior information at the design
stage, not about biological plausibility: the most faithful
generative model is not always the most estimable one, and when
independent evidence on the two channels is unavailable the
single-parameter pure architectures remain the more disciplined
default. Its appropriate
use cases are formal sensitivity analyses that ask how power
estimates change across the \((c_{bm,A}, c_{bm,B})\) grid when the
channel structure is genuinely uncertain, and evaluation studies
designed to determine which channel is dominant for a given
biomarker-drug pair.

In the prazosin-PTSD context, the dual-channel architecture captures the
possibility that systolic blood pressure simultaneously governs the
expected prazosin effect at the mean level (consistent with the
pharmacological argument that noradrenergic tone determines the
therapeutic window) and alters the coupling between blood pressure
and symptom change during active drug exposure (consistent with a
signaling-pathway interpretation). We regard the dual-channel architecture as the
appropriate framework for the prazosin-blood-pressure interaction if
and when mechanistic data become available to calibrate both
channels separately.

\end{orig}

\bigskip\noindent

\rule{\textwidth}{0.4pt}\bigskip

\section{6. Discussion}

\begin{rgt}
To be completed by rgt.

\end{rgt}

\begin{orig}
The combined architecture places the two architectures of the
companion paper within a single model, with each recovered as a
single-channel limit. This reframes the A-versus-B contrast as a
question about which channel weight is nonzero rather than a choice
between unrelated data-generating conventions. The super-additivity
result of Section\textasciitilde3 explains, rather than merely reports, the
compounding of the two channels: it follows from the additivity of
the biomarker-treatment signal and the convexity of the power
function, and is therefore a property of the estimand and the test,
not of the particular calibration. As in the companion paper, the
analysis model is held fixed; the conservatism of the standard
\texttt{corCAR1} Wald test, and the cluster-robust remedy for it,
are treated in the companion calibration study\textsuperscript{\citeproc{ref-pmsimstats-paper10}{5}} and
apply here as a common-mode effect across the channel-weight grid.
The combined form itself is not new to this series: a companion
paper\textsuperscript{\citeproc{ref-pmsimstats-paper03}{3}} introduces it conceptually as the general node
of a biomarker-moderation spectrum and supplies its psychometric
lineage\textsuperscript{\citeproc{ref-arminger1999}{2},\citeproc{ref-verbekelesaffre1996}{4}}. The contribution of
the present paper is complementary and specific: the analytical
super-additivity result of Section\textasciitilde3 and the calibrated
channel-weight simulation of Section\textasciitilde4, neither of which the
conceptual treatment provides.

\end{orig}

\bigskip\noindent

\rule{\textwidth}{0.4pt}\bigskip

\section{Conflict of interest}

The pmsimstats team declares no conflicts of interest.

\section{Funding}

This work received no specific funding.

\section{Reproducibility}

This manuscript was rendered with R\textasciitilde4.6.0 inside the project's
zzcollab Docker container (\texttt{rocker/tidyverse:4.6.0}; package
set captured in \texttt{renv.lock}). The combined-architecture
driver is at
\texttt{analysis/scripts/dgp-combined/01-run-combined-factorial.R}
and uses \texttt{furrr::future\_map\_dfr} with
\texttt{plan(multicore)}; it sets \texttt{set.seed(20260610)} and
\texttt{furrr\_options(seed = 20260610L)}. Per-replicate outputs are
at
\texttt{analysis/data/dgp-combined/01-combined-replicates.rds} and
cell-level summaries at \texttt{01-combined-summary.rds}. The matched
boundary re-run described in Section\textasciitilde4.2 will be added as a separate
driver that adopts the per-randomization-path sample-size convention
of the companion paper's production driver
(\texttt{analysis/scripts/quick-sim/01-dgp-prototype.R}). The
manuscript source, bibliography, and CSL file are at
\texttt{analysis/report/11-combined-dgp-architecture/}.

\section{References}\label{references}

\protect\phantomsection\label{refs}
\begin{CSLReferences}{0}{1}
\bibitem[\citeproctext]{ref-pmsimstats-paper01}
\CSLLeftMargin{1. }%
\CSLRightInline{{pmsimstats team}. Two architectures for simulating biomarker-treatment interactions: Implications for statistical power under carryover. Published online 2026.}

\bibitem[\citeproctext]{ref-arminger1999}
\CSLLeftMargin{2. }%
\CSLRightInline{Arminger G, Stein P, Wittenberg J. Mixtures of conditional mean- and covariance-structure models. \emph{Psychometrika}. 1999;64(4):475-494.}

\bibitem[\citeproctext]{ref-pmsimstats-paper03}
\CSLLeftMargin{3. }%
\CSLRightInline{{pmsimstats team}. Latent-class and mixture-model formulations for biomarker-treatment interaction in {N}-of-1 trials. Published online 2026.}

\bibitem[\citeproctext]{ref-verbekelesaffre1996}
\CSLLeftMargin{4. }%
\CSLRightInline{Verbeke G, Lesaffre E. A linear mixed-effects model with heterogeneity in the random-effects population. \emph{Journal of the American Statistical Association}. 1996;91(433):217-221.}

\bibitem[\citeproctext]{ref-pmsimstats-paper10}
\CSLLeftMargin{5. }%
\CSLRightInline{{pmsimstats team}. Type {I} calibration of fixed-effect interaction tests in linear mixed models: A staged diagnosis of working-covariance misspecification and a cluster-robust remedy. Published online 2026.}

\end{CSLReferences}

\end{document}
